\section{Branching, extinction and quasi-stationarity}
\label{sec:gw}

The linear birth--death process resolves events one at a time.  A
Galton--Watson process instead advances by generations, replacing every member of
the current population simultaneously by a random number of descendants.  This
change of clock preserves the branching property but alters the states available
to a surviving population.  That difference will eventually account for the gap
between \(A_{\mathrm c}(p)\) and its discrete counterpart.

\subsection{The Galton--Watson construction}

Let \(\{\xi_{n,i}:n\geq0,i\geq1\}\) be independent copies of a non-negative
integer-valued random variable \(\xi\), with
\[
  \Prb(\xi=k)=p_k,
  \qquad
  \phi(s)=\E[s^\xi]=\sum_{k=0}^{\infty}p_ks^k,
  \qquad
  m=\E[\xi]=\phi'(1).
\]
Starting from one ancestor, \(Z_0=1\), define
\begin{equation}
  Z_{n+1}=\sum_{i=1}^{Z_n}\xi_{n,i},
  \label{eq:gw-definition}
\end{equation}
with an empty sum equal to zero.  Conditional independence gives
\begin{equation}
  \E[s^{Z_{n+1}}\mid Z_n]=\phi(s)^{Z_n}.
  \label{eq:gw-conditional-pgf}
\end{equation}
Consequently the PGF \(F_n(s)=\E[s^{Z_n}]\) satisfies
\(F_{n+1}=F_n\circ\phi=\phi^{\circ(n+1)}\), while
\begin{equation}
  \E[Z_{n+1}\mid Z_n]=mZ_n,
  \qquad
  \E[Z_n]=m^n.
  \label{eq:gw-mean}
\end{equation}
The extinction probability by generation \(n\) is
\(u_n=F_n(0)=\Prb(Z_n=0)\).  Since \(F_{n+1}(0)=\phi(F_n(0))\),
\begin{equation}
  u_{n+1}=\phi(u_n),
  \qquad u_0=0.
  \label{iterPGF}
\end{equation}
The ultimate extinction probability is the smallest fixed point of \(\phi\) in
\([0,1]\) under the usual non-degeneracy assumptions
\cite{Harris1963,Athreya1972}.  In particular, extinction is certain when
\(m\leq1\), whereas \(m>1\) permits survival with positive probability.

\subsection{The multi-type extension}
\label{sec:multitype}

Later compartment models distinguish particles by type, location or phenotype.
The Galton--Watson construction extends without changing its logic.  Suppose
there are \(d\) types and a type-\(i\) individual produces the offspring vector
\(\boldsymbol\xi^{(i)}=(\xi^{(i)}_1,\ldots,\xi^{(i)}_d)\), with PGF
\begin{equation}
  \Phi_i(\boldsymbol s)
  =\E\!\left[\prod_{j=1}^{d}s_j^{\xi^{(i)}_j}\right],
  \qquad i=1,\ldots,d.
  \label{eq:multitype-offspring-pgf}
\end{equation}
The mean progeny matrix
\begin{equation}
  M_{ij}
  =\left.\frac{\partial\Phi_i}{\partial s_j}\right|_{\boldsymbol s=\boldsymbol1}
  =\E[\xi^{(i)}_j]
  \label{eq:mean-progeny-matrix}
\end{equation}
records the expected number of type-\(j\) children produced by one type-\(i\)
parent.  If population composition is represented by a row vector
\(\boldsymbol Z_n\), then
\begin{equation}
  \E[\boldsymbol Z_{n+1}\mid\boldsymbol Z_n]
  =\boldsymbol Z_nM,
  \qquad
  \E[\boldsymbol Z_n]=\E[\boldsymbol Z_0]M^n.
  \label{eq:multitype-mean}
\end{equation}

Let \(u_i^{(n)}\) be the probability that all descendants are extinct by
generation \(n\) when the founder has type \(i\).  Componentwise first-step
conditioning gives
\begin{equation}
  \boldsymbol u^{(n+1)}
  =\boldsymbol\Phi(\boldsymbol u^{(n)}),
  \qquad
  \boldsymbol u^{(0)}=\boldsymbol0,
  \label{eq:multitype-extinction-iteration}
\end{equation}
and the ultimate extinction vector is the least fixed point of
\(\boldsymbol\Phi\) in \([0,1]^d\).  For a finite, irreducible and
non-degenerate type system, the threshold is the spectral radius
\(\varrho(M)\): extinction is certain when \(\varrho(M)\leq1\), whereas
\(\varrho(M)>1\) gives a positive survival probability
\cite{Athreya1972}.  This criterion describes reproductive amplification of a
type mixture.  It does not determine rupture or killing when the catastrophe
rate depends on composition; that requires the transient kernel of the killed
process as well as its mean matrix.

\subsection{Binary death or division}
\label{sec:deathordivision}

The discrete model used throughout the chapter has the offspring law
\begin{equation}
  \Prb(\xi=2)=p,
  \qquad
  \Prb(\xi=0)=1-p,
  \qquad
  \phi(s)=1-p+ps^2.
  \label{eq:pgfbinary}
\end{equation}
Its mean offspring number is \(m=2p\), so the process is subcritical, critical or
supercritical according as \(p<\tfrac12\), \(p=\tfrac12\) or
\(p>\tfrac12\).  The mean \(\E[Z_n]=(2p)^n\) identifies these three regimes, but
it cannot distinguish a common early extinction from a rare surviving tree.

\begin{figure}[htbp]
  \centering
  \includegraphics[width=0.78\textwidth]{simpleGWvis.png}\\[5pt]
  \includegraphics[width=0.78\textwidth]{subGWvis.png}
  \caption{Genealogies of the binary Galton--Watson process, with generation
  increasing to the right.  The upper realisation is supercritical
  (\(p=0.58\)); the lower is subcritical (\(p=0.47\)).  A single trajectory does
  not identify the regime, but the contrast indicates why extinction and
  conditional size must be studied separately.}
  \label{fig:gwvis}
\end{figure}

Let
\begin{equation}
  S_n=\Prb(Z_n>0)=1-u_n
  \label{eq:Sn-def}
\end{equation}
denote survival to generation \(n\).  Substitution of
\(u_n=1-S_n\) into \eqref{iterPGF} yields the central recursion
\begin{equation}
  S_{n+1}=2pS_n-pS_n^2
  =mS_n\left(1-\frac{S_n}{2}\right),
  \qquad S_0=1.
  \label{SnIter}
\end{equation}
The quadratic term has a direct first-step interpretation.  If the founder dies,
the lineage is already extinct; if it divides, survival requires at least one of
the two independent daughter lineages to remain alive.

Because extinction is absorbing, \(S_n\) is non-increasing and has a limit.  The
fixed points of \eqref{SnIter} are
\begin{equation}
  S_\infty=0,
  \qquad
  S_\infty=\frac{2p-1}{p}.
  \label{eq:binary-fixed-points}
\end{equation}
The second value is admissible only when \(p>\tfrac12\).  Hence
\begin{equation}
  \Prb(Z_n>0\ \text{for all }n)
  =\begin{cases}
     0,&0\leq p\leq\tfrac12,\\[2pt]
     (2p-1)/p,&\tfrac12<p\leq1.
   \end{cases}
  \label{eq:binary-survival-limit}
\end{equation}
At \(p=\tfrac12\), the positive fixed point meets the origin and the linear
contraction disappears.  This is the source of the slow critical decay.

\subsection{Critical survival and lifetime}
\label{sec:criticalcase}

At criticality, \eqref{SnIter} becomes
\begin{equation}
  S_{n+1}=S_n-\frac12S_n^2.
  \label{eq:critical-survival-recurrence}
\end{equation}
The asymptotic form can be obtained without replacing the difference equation by
an ODE.  Since \(S_n\downarrow0\),
\begin{equation}
  \frac{1}{S_{n+1}}-\frac{1}{S_n}
  =\frac{1}{2-S_n}\longrightarrow\frac12.
  \label{eq:critical-reciprocal-increment}
\end{equation}
Ces\`aro averaging of these increments, or the Stolz--Ces\`aro theorem, therefore
gives
\begin{equation}
  S_n\sim\frac{2}{n}
  \qquad (n\to\infty).
  \label{twoovrn_sn}
\end{equation}

Let \(T=\inf\{n\geq0:Z_n=0\}\) be the extinction generation.  The tail-sum
formula gives
\begin{equation}
  \E[T]
  =\sum_{n=0}^{\infty}\Prb(T>n)
  =\sum_{n=0}^{\infty}S_n.
  \label{eq:ETsumSn}
\end{equation}
By \eqref{twoovrn_sn}, the last series has the same tail as twice the harmonic
series, and hence \(\E[T]=\infty\), although \(T<\infty\) almost surely.  The
apparently conflicting statements are resolved by conditioning:
\begin{equation}
  \E[Z_n\mid Z_n>0]=\frac{\E[Z_n]}{S_n}
  =\frac{1}{S_n}\sim\frac{n}{2}.
  \label{eq:critical-conditional-mean}
\end{equation}
Surviving critical trees become linearly larger, while their probability falls
like \(1/n\), leaving the unconditional mean equal to one.

Two related exponents follow from the same binary structure.  First,
\(
  \Prb(T=n)=S_{n-1}-S_n=\tfrac12S_{n-1}^2\sim2n^{-2}.
\)
Second, if
\(\mathcal T=\sum_{j\geq0}Z_j\) is total progeny, then a complete binary tree with
\(k\) division vertices has \(2k+1\) vertices altogether.  At \(p=\tfrac12\),
\begin{equation}
  \Prb(\mathcal T=2k+1)
  =\frac{1}{2^{2k+1}}\frac{1}{k+1}\binom{2k}{k}
  \sim \frac{1}{2\sqrt{\pi}}k^{-3/2}.
  \label{eq:critical-total-progeny}
\end{equation}
Thus the familiar \(3/2\) exponent applies along the odd support of
\(\mathcal T\), rather than to every integer indiscriminately.

\begin{figure}[htbp]
  \centering
  \includegraphics[width=0.96\textwidth]{power_law_fixed.pdf}
  \caption{Critical binary branching in \(10^6\) simulated realisations.  The
  panels show the \(n^{-1}\) survival tail, the \(n^{-2}\) extinction-time mass
  and the \(n^{-3/2}\) total-progeny mass.  The last scaling is understood on the
  odd support imposed by binary division.}
  \label{fig:power_law}
\end{figure}

\subsection{The limiting conditional distribution}
\label{sec:qsd}

Return to the subcritical range \(0<p<\tfrac12\), so \(m=2p<1\).  The
unconditional law converges to the point mass at zero.  Its informative remnant is
the conditional law
\begin{equation}
  \Prb(Z_n=k\mid Z_n>0),
  \qquad k\geq1.
  \label{eq:conditional-law}
\end{equation}
For a non-degenerate subcritical Galton--Watson process with the mild moment
conditions satisfied here, Yaglom's theorem states that
\eqref{eq:conditional-law} converges to a proper distribution on the positive
integers \cite{Yaglom1947,Athreya1972}.  This limiting conditional distribution
is quasi-stationary: if it is used as the initial law, one update followed by
conditioning on non-extinction leaves it unchanged.

The late survival scale is obtained directly from \eqref{SnIter}.  Since
\(S_n\to0\), the ratio \(S_{n+1}/S_n=m(1-S_n/2)\) tends to \(m\).  Define
\begin{equation}
  A(p)=\lim_{n\to\infty}\frac{S_n}{m^n},
  \qquad m=2p.
  \label{eq:Ap-def}
\end{equation}
The product argument in \cref{sec:product} proves that this limit exists and is
strictly positive for \(0<p<\tfrac12\).  It follows that
\begin{equation}
  S_n\sim A(p)m^n.
  \label{lateS}
\end{equation}
The quantity on the right of \eqref{eq:Ap-def} is a limit of ratios; it is not the
limit of \(S_n\), which remains zero.

The first conditional moment now follows without further approximation:
\begin{equation}
  \E[Z_n\mid Z_n>0]
  =\frac{\E[Z_n]}{S_n}
  =\frac{m^n}{S_n}
  \longrightarrow \frac{1}{A(p)}.
  \label{eq:qsmean}
\end{equation}
The constant \(1/A(p)\) will be called the quasi-stationary mean.  It measures
the size of a late surviving population, not the frequency with which survival
occurs.

The second moment makes the same separation visible.  The offspring variance is
\(
  \sigma^2=4p(1-p)
\), and the standard branching recurrence gives, for \(m\neq1\),
\begin{equation}
  \Var(Z_n)
  =\sigma^2m^{n-1}\frac{1-m^n}{1-m},
  \qquad
  \E[Z_n^2]=m^{2n}+\Var(Z_n).
  \label{eq:gw-second-moment}
\end{equation}
Dividing by \(S_n\sim A(p)m^n\) and using
\(
  \sigma^2/m=2(1-p)=2-m
\)
gives
\begin{align}
  \E[Z_n^2\mid Z_n>0]
  &\longrightarrow
  \frac{1}{A(p)}\left(1+\frac{1}{1-2p}\right),
  \label{eq:conditional-second-moment}\\
  \Var(Z_n\mid Z_n>0)
  &\longrightarrow
  \frac{1}{A(p)}\left(1+\frac{1}{1-2p}\right)
  -\frac{1}{A(p)^2}.
  \label{eq:condvar}
\end{align}
These moment limits are consequences of the Yaglom limit, rather than a proof of
distributional convergence by themselves.  They show that the whole late
conditional scale is controlled by \(p\) and the single constant \(A(p)\).

If \(H\) denotes the PGF of the Yaglom distribution, then the branching property
also yields the functional equation
\begin{equation}
  H(\phi(s))=mH(s)+1-m,
  \qquad H(0)=0,\quad H(1)=1.
  \label{eq:yaglom-functional}
\end{equation}
For the binary process its support lies in the positive even integers, because
\(Z_n\) is even for every \(n\geq1\).  This simple parity constraint will later
provide a useful bound on \(A(p)\).

\subsection{A continuum approximation and its limit}
\label{sec:riccati}

Write \(\varepsilon=1-2p=1-m>0\).  The exact survival recursion can be written as
\begin{equation}
  S_{n+1}-S_n=-\varepsilon S_n-pS_n^2.
  \label{eq:Sdiff}
\end{equation}
When \(S_n\) changes little over one generation, replacing the left-hand side by
a derivative gives the Riccati equation
\begin{equation}
  \frac{\dd S}{\dd n}=-\varepsilon S-pS^2.
  \label{eq:riccati}
\end{equation}
With \(U=1/S\), this becomes \(U'=\varepsilon U+p\), and hence
\begin{equation}
  S(n)=\frac{\varepsilon}
  {C\varepsilon\e^{\varepsilon n}-p}
  \label{eq:riccatisol}
\end{equation}
for a constant \(C\).  The approximation recovers the critical \(1/n\) scale and
the subcritical exponential scale, but it cannot determine their amplitude:
\(C\) is fixed by the early generations, exactly where the continuum replacement
is least accurate.  The constant \(A(p)\) therefore has to be recovered from the
discrete recursion itself.

\FloatBarrier
